\section{系统实现}

项目不仅实现了算法流程，还提供了面向使用者的多种调用方式。系统调用链路如图
\ref{fig:system-chain} 所示。网页、REST API、MCP 服务和命令行都可以触发同一
套编译流程，并返回路由后电路、SWAP 数、深度、耗时、轨迹和 SABRE 基线指标。

\begin{figure}[H]
  \centering
  \begin{tikzpicture}[
  node distance=0.75cm and 0.7cm,
  entry/.style={rectangle, rounded corners=2pt, draw=orange!65!black, fill=orange!8,
    align=center, minimum width=2.6cm, minimum height=0.78cm, font=\small},
  routeCore/.style={rectangle, rounded corners=2pt, draw=blue!60!black, fill=blue!6,
    align=center, minimum width=3.0cm, minimum height=0.85cm, font=\small},
  data/.style={rectangle, rounded corners=2pt, draw=green!45!black, fill=green!7,
    align=center, minimum width=3.0cm, minimum height=0.85cm, font=\small},
  arrow/.style={-{Latex[length=2mm]}, thick, draw=gray!70!black}
]
  \node[entry] (web) {网页演示};
  \node[entry, below=of web] (rest) {REST API};
  \node[entry, below=of rest] (mcp) {MCP 服务};
  \node[entry, below=of mcp] (cli) {命令行};

  \node[routeCore, right=2.0cm of rest] (request) {统一编译请求\\QASM 或示例};
  \node[routeCore, right=of request] (runtime) {NPQR 路由运行时};
  \node[data, above=of runtime] (model) {默认推理模型};
  \node[data, below=of runtime] (baseline) {SABRE 基线对比};
  \node[routeCore, right=of runtime] (result) {路由结果\\轨迹与指标};

  \draw[arrow] (web.east) -- (request.west);
  \draw[arrow] (rest.east) -- (request.west);
  \draw[arrow] (mcp.east) -- (request.west);
  \draw[arrow] (cli.east) -- (request.west);
  \draw[arrow] (request) -- (runtime);
  \draw[arrow] (model) -- (runtime);
  \draw[arrow] (runtime) -- (baseline);
  \draw[arrow] (runtime) -- (result);
  \draw[arrow] (baseline.east) -| (result.south);
\end{tikzpicture}

  \caption{系统调用链路}
  \label{fig:system-chain}
\end{figure}

\subsection{接口层}

REST API 适合网页和普通后端调用。用户可以提交内置示例名称或自定义 OpenQASM 2
文本，系统返回结构化 JSON 结果。MCP 服务适合工具客户端和自动化场景，能够查询
状态、列出示例、运行 NPQR、运行 SABRE 基线并返回算法证据。命令行适合本地检查、
演示和快速基线矩阵生成。

这些入口的共同点是都把 NPQR 作为默认编译流程，同时把 SABRE 作为对比结果返回。
接口中的对比字段用于解释基线表现，不表示 NPQR 失败时会由 SABRE 替代完成。

\subsection{编译层}

编译层负责把输入电路转化为路由任务，调用 NPQR 运行时，并将结果重新组织为用户
可读的输出。输出包括路由后 QASM、插入 SWAP 数、电路深度、运行时间和可复放轨迹。
轨迹字段使用户能够观察每一步 SWAP 和门执行，也便于报告中解释算法过程。

\subsection{模型与证据层}

项目包含一个默认推理模型。该模型在搜索循环中提供动作偏好，不单独承担完整路由
生成。系统还提供机器可读的算法证据，说明默认路由、算法组成、SABRE 基线、能力
边界和固定困难样例的闭环验证结果。该证据用于保证报告、网页和接口描述保持一致。
